% 从旧数学附录迁出的高次剩余说明；不新增算法。
\section{素数模高次剩余与单位根}\label{knowledge-higher-residue}
实现见第~\pageref{compact-KthResidue}~页；依赖本书离散对数、线性同余和快速幂。
设 $p$ 为素数，$g$ 是原根，$k>0$。
$a=0$ 时根只有零；否则写 $a=g^b$、$x=g^y$，方程转为
\[
ky\equiv b\pmod{p-1}.
\]
令 $d=\gcd(k,p-1)$。当 $d\nmid b$ 时无解；否则指数解为
$y=y_0+i(p-1)/d$，$0\le i<d$。
因此用三个数存全部根：
\[
\texttt{first}=g^{y_0},\qquad
\texttt{ratio}=g^{(p-1)/d},\qquad \texttt{count}=d.
\]
依次乘 ratio 即可枚举；不保证数值递增，也不保证 first 是最小根。
ratio 是阶为 $d$ 的单位根，因而枚举恰有 $d$ 个不同的根。
模 $p=2$ 同样有效；指数同余的模数是 1，唯一非零根为 1。

接口支持 $p\le10^{12}$，需要调用者提供已验证的素数与原根；
$k$ 支持完整 unsigned long long，先对 $p-1$ 取模后解同余。
$a$ 支持完整 signed long long，先正规化。
返回空 optional 表示无解，返回 \texttt{\{0,1,1\}} 表示唯一零根。
给定原根后的期望时间、空间为 $O(\sqrt p)$；枚举全部根另外需要 $O(d)$ 时间。
模乘须先扩展到 128 位，不可用 64 位乘积直接取模。

测试穷举小于 100 的每个素数、全部原根及小指数，
用任意精度模幂逐个枚举候选根并比较整个集合；
覆盖负数、64 位边界、模 2、无解与零根。
另检查 998244352 个根的紧凑表示，以及素数 999999999989 上的植根测试；
后一个素数与原根的条件由试除及乘法阶证书独立确认。
普通与 ASan/UBSan 均通过，尚无专门在线 AC。

这是 \href{https://oi-wiki.org/math/number-theory/residue/}{OI Wiki 高次剩余} 中
离散对数转线性同余方法的素数模实现。
素数幂及非互素被开方数见下一节；一般合数模数见后文；改良 Tonelli--Shanks 尚需另补，
不能把本节状态解读为该页面全部覆盖。

\clearpage
\section{素数幂模高次剩余与根的提升}\label{knowledge-prime-power-residue}
实现见第~\pageref{compact-PrimePowerRoots}~页。
令 $M=p^e$。当 $a=0$ 时，全部根是 $p^{\lceil e/k\rceil}$ 的倍数。
否则令 $t=v_p(a)<e$，必须有 $k\mid t$，否则无解。
取 $s=t/k$，$q=p^{e-t}$，写 $x=p^s y$，只需求单位方程
\[
y^k\equiv a/p^t\pmod q.
\]
每个单位根 $y$ 对应原模数下的根
\[
x=p^s y+\ell p^s q,\qquad 0\le\ell<p^{t-s}.
\]
这些提升保持互不相同，也解释了非互素被开方数可能产生大量根。

奇素数时，提供模 $p^e$ 的原根，约化后仍是模 $q$ 的原根，
沿用离散对数转线性同余的方法。
$q=2$ 时单位根只有 1；$q\ge4$ 且为二的幂时，单位表示为
$(-1)^z5^y$。符号指数模 2，5 的指数模 $q/4$；
$q=4$ 时只有符号分支。分别解两个指数同余即可，不可假设二的高次幂存在原根。

返回结构体保存 first、两个 ratio、两个 count，以及 step、lifts。
全部根由下式给出：
\[
\bigl(\texttt{first}\cdot\texttt{ratio}[0]^i\cdot\texttt{ratio}[1]^j\bmod\texttt{step}\bigr)
+\ell\cdot\texttt{step}.
\]
其中 $0\le i<\texttt{count}[0]$、$0\le j<\texttt{count}[1]$、
$0\le\ell<\texttt{lifts}$。
\texttt{get} 已处理模乘溢出，\texttt{size} 返回三者乘积；
不保证根按数值排序。零根情形复用同一表示。

模数上限 $10^{12}$，$k$ 支持完整无符号 64 位。
只检查范围，不检查 $p$ 的素性和奇素数原根的阶，调用者须保证。
默认参数 $g=0$ 仅用于 $p=2$；奇素数不能漏传原根。
时间和空间由 BSGS 决定，期望 $O(\sqrt q)$；
单次按索引取根使用快速幂，需 $O(\log M)$。

已穷举多个素数的全部不超过 256 的素数幂、所有被开方数与小指数，
用任意精度模幂直接枚举所有根，比较整个集合。
另测有符号与无符号 64 位边界、$2^{39}$ 的非平凡离散对数、
$3^{20}$ 的非互素方程，以及数千亿根的紧凑表示。
普通与 ASan/UBSan 均通过；尚无专门在线 AC。
一般合数的 CRT 组合见下一节；改良 Tonelli--Shanks 仍未完成，OI Wiki 页面继续标记部分覆盖。

\clearpage
\section{一般合数模高次剩余：CRT 组合}\label{knowledge-composite-residue}
实现见第~\pageref{compact-CompositeRoots}~页。
调用者提供完整分解 $M=\prod_{t=0}^{F-1}m_t$，其中 $m_t=p_t^{e_t}$ 两两互素，
并提供每个奇素数幂的原根。先分别求局部根族；任一局部无解，整体即无解。
若局部根数为 $c_t$，则整体根数为 $\prod_t c_t$。
只保存局部根族，不创建这个大小的数组。

对每个因子预处理
\[
E_t=\frac M{m_t}\left(\left(\frac M{m_t}\right)^{-1}\bmod m_t\right)\bmod M.
\]
对一组局部根 $r_t$，唯一对应的标准代表是
\[
x=\sum_t r_tE_t\bmod M.
\]
逆元存在是因为因子两两互素。
\texttt{get(index)} 将索引按各个 $c_t$ 转成混合进制数字，
再把局部索引拆成两个乘法根族索引和加法提升索引，最后用上述公式组合。
CRT 双射保证不重复、不遗漏，输出顺序不保证递增。

$M\le10^{12}$，$k>0$；空因子列表表示 $M=1$，返回唯一代表 0。
\texttt{mod} 保存模数，\texttt{total} 保存总解数，
\texttt{parts} 和 \texttt{weight} 保存局部根族与 $E_t$。
模数乘积检查上界，CRT 的乘积先扩展至 128 位。
单次按索引查询为 $O(F\log M)$；求解的主要时间和空间为局部 BSGS，
因子列表接口不包含整数分解与原根搜索的成本；自动准备重载见下一节。

已对 $1\le M\le180$、所有被开方数与小指数直接枚举所有根，
与两套代码的完整输出集合比较，并反转因子顺序重新检查。
另覆盖有符号与无符号 64 位边界、混合素数幂上的非互素方程和大量零根，
以及十一种不同素因子组成模数时的全部 1024 个平方根。
普通与 ASan/UBSan 均通过；尚无专门在线 AC。
至此已有已知分解及原根时的 BSGS 求根路线，
更快的改良 Tonelli--Shanks 等内容仍未覆盖，OI Wiki 条目保留部分完成状态。

\clearpage
\section{自动准备模数分解与原根}\label{knowledge-residue-setup}
实现见第~\pageref{compact-root_factors}~页。
\texttt{root\_factors(M,rho)} 把 $M\le10^{12}$ 分解为互异素数的幂，
返回合数求根需要的 $(p,e,g)$ 列表；模数 1 返回空列表。
二的幂不需要原根，固定填零。
奇素数幂 $m=p^e$ 的乘法群是循环群，设 $\varphi=m/p\cdot(p-1)$。
分解 $\varphi$ 后，候选 $g$ 为原根当且仅当
\[
p\nmid g,\qquad g^{\varphi/q}\not\equiv1\pmod m
\quad\text{对每个不同素因子 }q\mid\varphi.
\]
从 2 起枚举候选，返回最小原根；所有奇素数幂均有原根。
这里检查的是模 $p^e$ 的阶，不能仅取一个模 $p$ 的原根而不验证提升。

重载 \texttt{CompositeRoots::solve(a,k,M,rho)} 自动调用此准备过程。

随机分解工作区类型为 \texttt{PollardRho}。
如果需要在同一个模数下处理多个方程，先保存因子列表，调用原有列表接口，
避免重复分解和原根搜索。
准备成本包括 $M$ 与各 $\varphi$ 的因数分解，以及逐个候选的模幂检验；
不能把这部分直接算作常数时间。

独立试除验证 $1\ldots500$ 的全部模数与多组接近 $10^{12}$ 的边界值，
并用任意精度模幂验证返回原根的阶。
小模数进一步检查更小候选不能生成全部单位。
直接传模数的两套求根接口，在小模数上与全部候选根的枚举结果比较。
普通与 ASan/UBSan 均通过；尚无自动准备入口的在线 AC。

\clearpage
